\documentclass[11pt]{article}
\usepackage{preamble}
\renewcommand{\answer}[1]{}

\begin{document}

\ladderheading{AP Statistics \textperiodcentered\ Unit 4 \textperiodcentered\ Topic 4.4 \textperiodcentered\ B: 2027-01-19 \textperiodcentered\ E: 2027-03-05}{Two Measurements per Technician}

\focusbanner{TODAY'S PROBLEM}

Suppose a company takes an SRS of 18 of its 300 technicians. Each completes the same task with old and new interfaces; interface order is randomly assigned. The plot shows $d=\text{new time}-\text{old time}$. Researchers ask whether the new interface reduces the population mean time.\par\smallskip \textbf{Nia:} ``The technician seems to be the independent unit, so I would build one difference per person; the subtraction order matters for the tail.''\par \textbf{Colin:} ``There are 36 recorded times, so I would use two independent samples of 18 and a two-sided alternative.''\par
\begin{center}
\begin{tikzpicture}
\begin{axis}[
  width=0.77\linewidth,
  height=1.44in,
  ymin=0, ymax=2, ytick=\empty, axis y line=none, axis x line=bottom,
  xlabel={New-interface time minus old-interface time (seconds)}, tick label style={font=\small}, label style={font=\small}
]
\addplot[only marks, mark=*, mark size=2.2pt, color=dayblue] coordinates {
  (-4.2,1)
  (-3.6,1)
  (-3.2,1)
  (-2.9,1)
  (-2.6,1)
  (-2.4,1)
  (-2.2,1)
  (-2.0,1)
  (-1.9,1)
  (-1.7,1)
  (-1.6,1)
  (-1.4,1)
  (-1.2,1)
  (-1.0,1)
  (-0.7,1)
  (-0.4,1)
  (0.0,1)
  (0.6,1)
};
\end{axis}
\end{tikzpicture}
\end{center}
\begin{firsttakebox}{FIRST TAKE --- before the video (5 min)}
Would you compare each technician's two times or treat them as separate people? Explain in one sentence.
\workspace{0.9in}
\end{firsttakebox}

\begin{callout}{\IconBook\ SET UP A TEST FOR A MEAN (keep on the page)}{calloutgreen}
For a population mean with unknown population standard deviation, use a one-sample $t$-test.
Matched pairs are two linked measurements: subtract in the same order and analyze one sample of differences.
Write hypotheses about the population mean, not the sample mean. The null uses equality; the alternative
uses $<$, $>$, or $\neq$, as determined by the research question and subtraction order.
Check random sampling or assignment, $n\leq0.10N$ when sampling without replacement, and shape.
For $n<30$, the values being analyzed should have no strong skewness or outliers.
\end{callout}

\newpage
\textbf{Finish after the video:}\par\smallskip

\dokbadge{1}~\textbf{(a)} Define $\mu_d$ in context.\par
\workspace{0.8in}

\dokbadge{2}~\textbf{(b)} Check the random-sample, 10 percent, and small-sample shape conditions in three short notes.\par
\workspace{1.4in}

\dokbadge{3}~\textbf{(c)} In 3--4 sentences, choose a setup and justify it. State the hypotheses using the given subtraction order and reduction question, then explain why the 36 times are not 36 independent people.\par
\workspace{2.0in}

\begin{sentenceframebox}\raggedright\textbf{Frame:} The independent units are \blank[1.4in]; preserving each pair changes the analysis by \blank[2.6in].\end{sentenceframebox}

\begin{sentenceframebox}\raggedright\textbf{Frame:} With the stated subtraction order, the claim corresponds to \blank[1.2in]; the rejected setup instead answers \blank[2.5in].\end{sentenceframebox}

\begin{summaryexitbox}{EXIT --- turn this in}
One thing the video changed about my first take: \hrulefill\par
\workspace{0.4in}
\end{summaryexitbox}

\end{document}
